\documentclass[a4paper,11pt]{article}

% =========================================================
% Standalone-Artikel (v2): Quadratsummen-Embedding von Primquadraten
% mit Lösungsraum, Symmetrien und Testprotokoll
% Datum: \today
% =========================================================

% --- Sprache / Kodierung ---
\usepackage[utf8]{inputenc}
\usepackage[T1]{fontenc}
\usepackage[ngerman]{babel}

% --- Mathe ---
\usepackage{amsmath,amssymb,amsthm}

% --- Layout / Links ---
\usepackage{geometry}
\geometry{left=25mm,right=25mm,top=25mm,bottom=25mm}
\usepackage{hyperref}
\hypersetup{
	colorlinks=true,
	linkcolor=blue,
	urlcolor=blue,
	citecolor=blue
}

% --- Grafiken (robust; safe include) ---
\usepackage{graphicx}
\DeclareGraphicsExtensions{.pdf,.png,.jpg}
\graphicspath{{./}{fig/}{figs/}{images/}{img/}}
\newcommand{\includegraphicssafe}[2][]{%
	\IfFileExists{#2}{\includegraphics[#1]{#2}}{%
		\fbox{\parbox{0.95\linewidth}{\textbf{Missing graphic:} \texttt{#2}}}%
	}%
}

% --- Theorem-Umgebungen (deutsch) ---
\newtheorem{satz}{Satz}[section]
\newtheorem{lemma}[satz]{Lemma}
\newtheorem{proposition}[satz]{Proposition}
\newtheorem{korollar}[satz]{Korollar}
\theoremstyle{definition}
\newtheorem{definition}[satz]{Definition}
\theoremstyle{remark}
\newtheorem{bemerkung}[satz]{Bemerkung}
\newtheorem{beispiel}[satz]{Beispiel}

% --- Makros ---
\newcommand{\Z}{\mathbb{Z}}
\newcommand{\N}{\mathbb{N}}
\newcommand{\C}{\mathbb{C}}
\newcommand{\Zi}{\mathbb{Z}[i]}
\newcommand{\Norm}{\mathrm{N}}
\newcommand{\Arg}{\mathrm{arg}}
\newcommand{\Ree}{\mathrm{Re}}
\newcommand{\Imm}{\mathrm{Im}}
\newcommand{\abs}[1]{\left|#1\right|}
\newcommand{\unit}{\{\pm 1,\pm i\}}

% --- Titel ---
\title{\textbf{Quadratsummen-Embedding von Primquadraten in \(\Zi\):\\
		Lösungsraum, Symmetrien, Korridorwahl und prüfbare Spiraltests}}
\author{Thomas Hoffbauer}
\date{\today}

\begin{document}
	\maketitle
	
	\begin{abstract}
		Wir formulieren ein arithmetisch-geometrisches Framework: Primquadrate \(p^2\) werden als Normen
		\(p^2=\Norm(z)\) von Gaußzahlen \(z\in\Zi\) interpretiert und damit als Gitterpunkte \((x,y)\in\Z^2\) auf Kreisen
		\(|z|=p\) in der komplexen Ebene dargestellt. Ein ``Korridor'' um die imaginäre Achse wird als Auswahlregel (Gauge-Fix)
		formalisiert. Zentral ist die klare Trennung zwischen (i) klassischer Zahlentheorie (Restklassen, Summe-zweier-Quadrate,
		Zerlegung in \(\Zi\)) und (ii) modellhaften Auswertungen (achsennahe Selektion, Spiral-Metriken). Wir ergänzen eine
		systematische Beschreibung des Lösungsraums von \(x^2+y^2=p^2\) bis auf Einheiten sowie ein explizites,
		reproduzierbares Testprotokoll (Nullhypothese, Robustheit bzgl.\ Auswahlregeln, Unwrapped-Angle).
	\end{abstract}
	
	\tableofcontents
	
	% =========================================================
	\section{Einleitung: Normgeometrie statt Behauptungsgeometrie}
	% =========================================================
	Eine Darstellung
	\[
	n=x^2+y^2
	\]
	entspricht in \(\Zi\) einer Normdarstellung
	\[
	n=\Norm(x+iy)=(x+iy)(x-iy).
	\]
	Damit werden Zahlen zu geometrischen Objekten: \(n\) ist Norm eines Gitterpunktes, also Radiusquadrat eines Kreises
	im \(\Z^2\)-Gitter.
	
	Ziel dieses Textes ist \emph{nicht} die Behauptung einer ``tiefen Spiralwahrheit'' über Primzahlen, sondern:
	\begin{itemize}
		\item ein wohldefiniertes \emph{Embedding} \(p\mapsto z_p\in\Zi\) mit \(\Norm(z_p)=p^2\),
		\item eine wohldefinierte \emph{Korridor-/Gauge-Fix}-Regel zur Auswahl eines Repräsentanten pro Symmetrieorbit,
		\item und ein \emph{prüfbares Testprotokoll}, ob die daraus entstehenden Winkel \(\theta_p=\Arg(z_p)\) in geeigneten
		Koordinaten (z.\,B.\ \(\log p\) gegen unwrapped \(\theta_p\)) Spiralregularitäten zeigen.
	\end{itemize}
	
	% =========================================================
	\section{Primklassen \(\bmod 12\) und Quadrat-Kollaps nach \(1\)}
	% =========================================================
	\begin{definition}[Primklassen \(\bmod 12\)]
		Für jede Primzahl \(p>3\) gilt \(p\equiv 1,5,7,11\pmod{12}\).
		Wir definieren
		\[
		E:\ p\equiv 1\!\!\pmod{12},\quad
		a:\ p\equiv 5\!\!\pmod{12},\quad
		b:\ p\equiv 7\!\!\pmod{12},\quad
		c:\ p\equiv 11\!\!\pmod{12}.
		\]
	\end{definition}
	
	\begin{lemma}[Quadrat-Kollaps]
		\label{lem:kollaps}
		Für jede Primzahl \(p>3\) gilt \(p^2\equiv 1\pmod{12}\).
	\end{lemma}
	
	\begin{proof}
		Für \(p>3\) ist \(\gcd(p,12)=1\). Dann gilt \(p^2\equiv 1\pmod{3}\) und \(p^2\equiv 1\pmod{4}\), also \(p^2\equiv 1\pmod{12}\).
	\end{proof}
	
	\begin{bemerkung}
		Lemma~\ref{lem:kollaps} sagt nur etwas über Restklassen. Die ``Quadratsummen-Geometrie'' entsteht erst über
		Normdarstellungen in \(\Zi\).
	\end{bemerkung}
	
	% =========================================================
	\section{Quadratsummen in \(\Zi\): Existenz, Nichttrivialität, Zerlegung}
	% =========================================================
	\begin{definition}[Gauß-Norm]
		Für \(z=x+iy\in\Zi\) definieren wir \(\Norm(z)=x^2+y^2\).
	\end{definition}
	
	\begin{satz}[Summe zweier Quadrate]
		\label{thm:sos}
		Eine natürliche Zahl \(n\) ist genau dann Summe zweier Quadrate, wenn in der Primfaktorzerlegung von \(n\)
		jede Primzahl \(q\equiv 3\pmod 4\) mit \emph{geradem} Exponenten auftritt.
	\end{satz}
	
	\begin{korollar}[Primfälle und Primquadrate]
		\label{cor:primcases}
		Sei \(p\) eine Primzahl.
		\begin{enumerate}
			\item Falls \(p\equiv 1\pmod4\), existieren \(u,v\in\Z\) mit \(uv\neq 0\) und \(p=u^2+v^2\).
			\item Falls \(p\equiv 3\pmod4\), ist \(p\) keine nichttriviale Quadratsumme, aber \(p^2\) ist Quadratsumme.
			\item Für jedes \(p\) existiert mindestens eine Darstellung \(p^2=x^2+y^2\) (notfalls trivial \(p^2=p^2+0^2\)).
		\end{enumerate}
	\end{korollar}
	
	\begin{bemerkung}[Wichtig: keine automatische Achsennähe]
		Aus Korollar~\ref{cor:primcases} folgt \emph{nicht}, dass es für viele \(p\) nichttriviale Darstellungen von \(p^2\)
		mit \(\abs{x}/p\) sehr klein gibt. Eine ``achsennahe'' Lösung ist ein \emph{Selektionskriterium} und wird empirisch
		über Häufigkeiten untersucht.
	\end{bemerkung}
	
	% =========================================================
	\section{Lösungsraum von \(x^2+y^2=p^2\): primitive Lösungen und Symmetrieorbits}
	% =========================================================
	\begin{definition}[Primitive Lösung]
		Eine Lösung \((x,y)\in\Z^2\) von \(x^2+y^2=p^2\) heißt \emph{primitiv}, falls \(\gcd(x,y)=1\).
	\end{definition}
	
	\begin{lemma}[Primitivität bei Primquadraten]
		\label{lem:primitive}
		Sei \(p\) prim. Dann ist jede nichttriviale Lösung \(x^2+y^2=p^2\) mit \(xy\neq 0\) \emph{nicht} primitiv:
		Es gilt \(\gcd(x,y)=p\) oder \(\gcd(x,y)=1\) ist nur möglich, wenn \(p=1\) (nicht prim).
		Insbesondere sind die ``natürlichen'' nichttrivialen Lösungen typischerweise Vielfache eines primitiven Paares aus \(p=u^2+v^2\).
	\end{lemma}
	
	\begin{proof}
		Aus \(x^2+y^2=p^2\) folgt, dass jeder gemeinsame Teiler von \(x\) und \(y\) auch \(p\) teilt. Da \(p\) prim ist, ist
		\(\gcd(x,y)\in\{1,p\}\). Für \(xy\neq 0\) und \(p\) prim zeigt eine Standardargumentation über Kongruenzen mod \(p\),
		dass \(\gcd(x,y)=p\) in den nichttrivial konstruierten Fällen (z.\,B. aus \(p=u^2+v^2\) via Quadratur) tatsächlich eintritt.
	\end{proof}
	
	\begin{bemerkung}
		Die Geometrie lebt also nicht von ``primitiven Punkten'' auf dem Kreis \(|z|=p\), sondern von der kontrollierten
		Erzeugung nichttrivialer Punkte als strukturierte Vielfache (insbesondere über \(z_p=(u+iv)^2\)).
	\end{bemerkung}
	
	\begin{proposition}[Klassifikation der wesentlichen Kandidaten für \(p\equiv 1\pmod4\)]
		\label{prop:candidates}
		Sei \(p\equiv 1\pmod4\) prim und \(p=u^2+v^2\) (mit \(u,v>0\)). Setze \(\pi=u+iv\in\Zi\).
		Dann gilt \(p=\pi\overline{\pi}\) und \(p^2\) hat in \(\Zi\) (bis auf Einheiten und Konjugation) zwei natürliche
		Faktorformen:
		\[
		p^2=(\pi^2)(\overline{\pi}^2)\quad\text{und}\quad p^2=(p)(p).
		\]
		Die daraus induzierten Normdarstellungen in \(\Z^2\) liefern als ``wesentliche'' nichttriviale Kandidaten
		(unter Symmetrien \(\pm\) und Konjugation) genau die Punkte
		\[
		\pi^2=(u^2-v^2)+i(2uv),\qquad i\pi^2=-(2uv)+i(u^2-v^2),
		\]
		sowie deren Vorzeichen-/Konjugationsbilder. Die trivialen Achsenpunkte entsprechen \(ip\) (bzw. \(p\)).
	\end{proposition}
	
	\begin{proof}[Begründung (skizzenhaft, aber präzise)]
		Die Einheiten in \(\Zi\) sind \(\unit\). Multiplikation mit Einheiten rotiert um \(\pi/2\) und ändert die Norm nicht.
		Konjugation spiegelt an der reellen Achse. Für \(p\equiv 1\pmod4\) zerfällt \(p\) in \(\Zi\) als \(\pi\overline{\pi}\).
		Damit entsteht \(p^2\) als \((\pi^2)(\overline{\pi}^2)\). Das Element \(\pi^2\) liefert eine konkrete nichttriviale
		Darstellung von \(p^2\) als Summe zweier Quadrate (Quadratur-Identität). Rotationen durch Einheiten geben die
		weiteren Kandidaten. Die alternative Faktorisierung \(p^2=p\cdot p\) ist arithmetisch trivial und liefert die Achsenlösungen.
	\end{proof}
	
	\begin{bemerkung}[Warum das wichtig ist]
		Das beseitigt eine verbreitete Unschärfe: Für \(p^2\) gibt es nicht ``beliebig viele'' geometrisch unterschiedliche
		Darstellungen, sondern unter den natürlichen, strukturerhaltenden Kandidaten im Wesentlichen eine kleine, klar
		beschreibbare Familie (bis auf Symmetrien).
	\end{bemerkung}
	
	% =========================================================
	\section{Korridorwahl als Gauge-Fix: Definitionen ohne Overclaim}
	% =========================================================
	\begin{definition}[\(\varepsilon\)-Korridor um die imaginäre Achse]
		Fixiere \(\varepsilon\in(0,1)\). Ein Punkt \(z=x+iy\) mit \(|z|=p\) heißt \(\varepsilon\)-achsennahe, falls
		\[
		\abs{x}\le \varepsilon p.
		\]
	\end{definition}
	
	\begin{definition}[Quadratsummen-Embedding]
		Eine Abbildung \(\Phi\) ordnet jeder Primzahl \(p>3\) eine Gaußzahl \(z_p\in\Zi\) zu mit
		\[
		\Norm(z_p)=p^2.
		\]
		Eine \emph{korridorbasierte} Regel wählt \(z_p\) so, dass \(\abs{\Ree(z_p)}\) möglichst klein ist
		(z.\,B.\ in einem \(\varepsilon\)-Korridor), falls unter den ``wesentlichen Kandidaten'' (Proposition~\ref{prop:candidates})
		eine nichttriviale Option verfügbar ist.
	\end{definition}
	
	\begin{bemerkung}[Variationsprinzip statt Ästhetik]
		Die Achse ist keine ``Wahrheit'' über Primzahlen, sondern eine \emph{Modellentscheidung}:
		Man fixiert eine Referenzrichtung (Gauge) und minimiert die Abweichung davon. Das ist methodisch analog zur Wahl einer
		Repräsentantenkonvention in einem Symmetrieorbit.
	\end{bemerkung}
	
	% =========================================================
	\section{Kanonische Regeln R1--R3: was sie garantieren (und was nicht)}
	% =========================================================
	\begin{definition}[R1: Minimaler Realteil]
		Wähle \(z_p=x+iy\) unter allen (zugelassenen) Kandidaten so, dass \(\abs{x}\) minimal ist; bei Gleichstand wähle \(y>0\).
	\end{definition}
	
	\begin{definition}[R2: Minimaler Winkelabstand zur imaginären Achse]
		Wähle \(z_p\) so, dass \(\abs{\Arg(z_p)-\pi/2}\) minimal ist, unter \(\Imm(z_p)>0\).
	\end{definition}
	
	\begin{definition}[R3: Strukturregel über \(\pi^2\)]
		Falls \(p\equiv 1\pmod4\): bestimme \(p=u^2+v^2\), setze \(\pi=u+iv\) und wähle \(z_p=\pi^2\) oder \(i\pi^2\) nach
		Achsennähe (minimiere \(\abs{\Ree(z)}\)). Falls \(p\equiv 3\pmod4\): setze \(z_p=ip\) (trivial, achsenliegend).
	\end{definition}
	
	\begin{proposition}[R1/R2 sind Optimierer, keine Garantien]
		R1 minimiert \(\abs{\Ree(z_p)}\) über die zugelassene Kandidatenmenge; R2 minimiert den Winkelabstand zur Achse.
		Keines von beiden garantiert eine uniforme Schranke \(\abs{\Ree(z_p)}\le \varepsilon p\) für festes \(\varepsilon\),
		sondern definiert lediglich eine wohldefinierte Auswahl und damit eine testbare Statistik der Achsennähe.
	\end{proposition}
	
	% =========================================================
	\section{Spiraltests: definierte Messgrößen, Nullhypothese, Robustheit}
	% =========================================================
	\subsection{Daten und Koordinaten}
	\begin{definition}[Winkel- und Logradialdaten]
		Für \(z_p\neq 0\):
		\[
		r_p:=|z_p|=p,\qquad \theta_p:=\Arg(z_p)\in(-\pi,\pi],\qquad \rho_p:=\log p.
		\]
	\end{definition}
	
	\subsection{Unwrapped Angle}
	Die Periodizität von \(\theta\) macht lineare Fits mehrdeutig. Deshalb definieren wir eine ``entwrapte'' Winkelspur.
	
	\begin{definition}[Unwrapped Angle \(\widetilde{\theta}_p\)]
		Sortiere Primzahlen \(p_1<p_2<\dots<p_n\). Setze \(\widetilde{\theta}_{p_1}:=\theta_{p_1}\).
		Für \(k\ge 2\) wähle \(\widetilde{\theta}_{p_k}\) als den Wert der Form \(\theta_{p_k}+2\pi m\) mit \(m\in\Z\),
		der \(\abs{\widetilde{\theta}_{p_k}-\widetilde{\theta}_{p_{k-1}}}\) minimiert.
	\end{definition}
	
	\begin{bemerkung}
		Das ist eine wohldefinierte, implementierbare Konvention. Alternative Unwrap-Regeln (z.\,B.\ mit Glättung oder
		Sprungschranken) können als Robustheitscheck eingesetzt werden.
	\end{bemerkung}
	
	\subsection{Fit-Score und Korrelation}
	\begin{definition}[Lineares Spiralmodell und Score]
		Für Datenpaare \((\widetilde{\theta}_{p_k},\rho_{p_k})\) definiere den kleinsten Quadrate-Score
		\[
		S(\Phi;X):=\min_{A,B}\ \sum_{p_k\le X}\big(\rho_{p_k}-(A\widetilde{\theta}_{p_k}+B)\big)^2.
		\]
		Optional normiert man \(S\) durch die Varianz von \(\rho\) oder durch \(n\).
	\end{definition}
	
	\begin{definition}[Nullhypothese und Permutationstest]
		Als Nullhypothese \(H_0\) dient z.\,B.\ ``keine strukturelle Kopplung zwischen \(\theta_p\) und \(\log p\)''.
		Operationalisierung: Permutiere die Winkel \(\theta_{p_k}\) zufällig über die Primindices (oder ersetze sie durch i.i.d.\ Uniformwinkel),
		unwrappe erneut und berechne \(S\). Der p-Wert ist der Anteil der Nullmodelle mit \(S_{\mathrm{null}}\le S_{\mathrm{obs}}\).
	\end{definition}
	
	\begin{bemerkung}[Robustheit ist Pflicht]
		Der Score ist nur dann interessant, wenn er stabil bleibt gegen:
		(i) Wechsel der Auswahlregel (R1/R2/R3), (ii) Variation von \(\varepsilon\) bzw.\ Kandidatenmenge,
		(iii) alternative Unwrap-Konventionen.
	\end{bemerkung}
	
	\subsection{Häufigkeitsfunktion der Achsennähe}
	\begin{definition}[Achsennähe-Häufigkeit]
		Für \(\varepsilon\in(0,1)\) und Obergrenze \(X\) definieren wir
		\[
		f(\varepsilon;X):=\frac{1}{\pi(X)}\#\Big\{p\le X:\ \frac{\abs{\Ree(z_p)}}{p}\le \varepsilon\Big\}.
		\]
		Diese Funktion misst, wie oft die gewählte Regel \(\Phi\) tatsächlich in einem \(\varepsilon\)-Korridor landet.
	\end{definition}
	
	\begin{bemerkung}
		Das ist der sauberste Ersatz für implizite Aussagen wie ``entlang eines Korridors existieren Spiralen''.
		Hier wird klar: Der Korridor ist erst einmal eine \emph{Messgröße}.
	\end{bemerkung}
	
	% =========================================================
	\section{Kurzer Ausblick: Multiplikative Kompatibilität (optional)}
	% =========================================================
	\begin{bemerkung}[Homomorphie-Nähe ist nicht selbstverständlich]
		Ein verlockendes Ziel wäre eine Regel \(\Phi\) mit
		\[
		z_{pq}\approx z_p z_q
		\]
		(im Sinne von Winkeladditivität oder kleiner Korrekturen). Für kanonische Auswahlregeln ist dies typischerweise
		\emph{nicht} exakt, weil die Gauge-Fixierung nicht multiplikativ respektiert werden muss.
		Gerade diese Spannung kann jedoch als diagnostisches Werkzeug dienen: Wie stark ``zerreißt'' die Auswahlregel
		multiplikative Struktur?
	\end{bemerkung}
	
	% =========================================================
	% Upgrade-Baustein 1: Invarianztests (Rotation / Gauge / Robustheit)
	% =========================================================
	
	\section{Invarianztests: Rotation, alternative Gauges und Robustheitsprotokoll}
	\label{sec:invarianz}
	
	\subsection{Rotierte Referenzachse als systematischer Test}
	Die Korridorwahl um die imaginäre Achse ist eine Modellentscheidung (Gauge-Fix). Ein zentraler Test besteht darin,
	die Referenzachse zu rotieren und alle Kennzahlen erneut zu messen.
	
	\begin{definition}[Rotationsoperator in \(\Zi\)]
		Für \(\varphi\in[0,2\pi)\) definieren wir die Rotation auf \(\C\) durch
		\[
		R_\varphi(z):=e^{-i\varphi}z.
		\]
		Wir rotieren \emph{nicht} das Gitter, sondern die Koordinatenbeschreibung: Achsennähe relativ zu einer Achse bei Winkel \(\varphi\)
		ist Achsennähe relativ zur imaginären Achse nach Anwendung von \(R_\varphi\).
	\end{definition}
	
	\begin{definition}[\(\varepsilon\)-Korridor relativ zu Winkel \(\varphi\)]
		Ein Punkt \(z\in\C\) mit \(|z|=p\) liegt im \(\varepsilon\)-Korridor um die Achse mit Richtung \(\varphi\), wenn
		\[
		\left|\Re\left(R_\varphi(z)\right)\right|\le \varepsilon p.
		\]
		Für \(\varphi=0\) ist dies genau der ursprüngliche Korridor um die imaginäre Achse.
	\end{definition}
	
	\begin{bemerkung}
		Diese Definition macht den Rotationstest präzise: Jede ``Spiralstruktur'' oder ``Achsennähe'' darf nicht nur
		für einen einzelnen Winkel \(\varphi=0\) gelten, sondern muss im Idealfall \emph{stabil} (oder zumindest strukturiert nachvollziehbar)
		als Funktion von \(\varphi\) auftreten.
	\end{bemerkung}
	
	\subsection{Rotationstest-Protokoll (Sweep)}
	\begin{definition}[Rotationstest-Sweep]
		Fixiere einen Satz von Winkeln \(\Phi_\mathrm{rot}=\{\varphi_j\}_{j=0}^{m-1}\) mit \(\varphi_j=2\pi j/m\).
		Für jede Regel \(\Phi\) (R1/R2/R3 etc.) und jedes \(X\) berechne die Kennzahlen
		\[
		f(\varepsilon;X,\varphi_j),\qquad S(\Phi;X,\varphi_j),
		\]
		wobei \(f\) und \(S\) mit \(R_{\varphi_j}(z_p)\) statt \(z_p\) definiert werden.
	\end{definition}
	
	\begin{definition}[Rotations-Report (Minimalset)]
		Für jede Kombination \((\Phi,X)\) soll mindestens berichtet werden:
		\begin{enumerate}
			\item \(\varphi\mapsto f(\varepsilon;X,\varphi)\) für ein kleines Set \(\varepsilon\in\{\varepsilon_1,\varepsilon_2,\varepsilon_3\}\) (z.\,B.\ \(0.02,0.05,0.10\)).
			\item \(\varphi\mapsto S(\Phi;X,\varphi)\) (ggf. normalisiert).
			\item Ein Nullmodell-Band für \(S\) (z.\,B.\ 5\%–95\% Quantile aus Permutation/Uniform-Winkeln) pro \(\varphi\).
		\end{enumerate}
	\end{definition}
	
	\subsection{Alternative Gauge-Fixes: Kandidatenmengen und Auswahlregeln}
	Die stärkste Robustheit entsteht nicht nur durch Rotation, sondern auch durch Variation der \emph{Gauge-Fix-Regel}
	und der \emph{zugelassenen Kandidatenmenge}.
	
	\begin{definition}[Kandidatenmengen \(\mathcal{C}_p\)]
		Für \(p\equiv 1\pmod 4\) sei \(\pi_p\) eine kanonische Wahl mit \(p=\pi_p\overline{\pi_p}\). Dann definieren wir z.\,B.:
		\begin{align*}
			\mathcal{C}_p^{(1)} &:= \{\pi_p^2\},\\
			\mathcal{C}_p^{(2)} &:= \{\pi_p^2,\, i\pi_p^2\},\\
			\mathcal{C}_p^{(4)} &:= \{u\pi_p^2:\ u\in\{\pm 1,\pm i\}\},
		\end{align*}
		und optional auch das Hinzufügen der Konjugierten (Spiegelung). Für \(p\equiv 3\pmod 4\) setze \(\mathcal{C}_p:=\{ip\}\) als Basisfall.
	\end{definition}
	
	\begin{definition}[Gauge-Fix als Minimierungsregel]
		Eine Gauge-Fix-Regel wählt aus \(\mathcal{C}_p\) einen Repräsentanten
		\[
		z_p \in \arg\min_{z\in\mathcal{C}_p}\left|\Re\!\left(R_\varphi(z)\right)\right|,
		\quad \text{mit Tie-Breaker } \Im\!\left(R_\varphi(z)\right)>0.
		\]
		Damit sind R1/R2/R3 als Spezialfälle (bei geeigneter \(\mathcal{C}_p\)) einheitlich formulierbar.
	\end{definition}
	
	\subsection{Robustheitskriterium (konservativ)}
	\begin{definition}[Robustheitskriterium]
		Ein Effekt (z.\,B.\ ``Score kleiner als Nullmodell'' oder ``hohes \(f(\varepsilon)\)'') gilt nur dann als robust bis \(X\),
		wenn er für
		\begin{enumerate}
			\item mindestens zwei Kandidatenmengen \(\mathcal{C}^{(\cdot)}\),
			\item mindestens zwei Auswahlregeln (z.\,B.\ Minimierung von \(|\Re|\) vs. Minimierung von Winkelabstand),
			\item und für einen nichttrivialen Anteil der Rotationswinkel \(\varphi_j\)
		\end{enumerate}
		gleichzeitig auftritt.
	\end{definition}
	
	\begin{bemerkung}
		Dieses Kriterium ist absichtlich streng: Es verhindert, dass ein Artefakt einer einzelnen Konvention als ``Struktur'' verkauft wird.
	\end{bemerkung}
	
	
	% =========================================================
	% Upgrade-Baustein 2: Fourier-Momente und kanonische Wahl von \pi_p
	% =========================================================
	
	\section{Fourier-Momente der Winkel: kanonische \texorpdfstring{$\pi_p$}{pi\_p}-Wahl, Nullmodell und Auswertung}
	\label{sec:fourier}
	
	\subsection{Warum Fourier-Momente?}
	Wenn die Winkel der gaußschen Primfaktoren \(\pi_p\) in Sektoren ``gleichverteilt'' sind, sollten nichttriviale
	Fouriermomente stark auslöschen. Das ist die analytisch anschlussfähigste Messgröße neben dem Spiral-Fit.
	
	\subsection{Kanonische Wahl von \(\pi_p\) (Einheitenproblem lösen)}
	Für \(p\equiv 1\pmod 4\) zerfällt \(p\) in \(\Zi\) als \(p=\pi\overline{\pi}\), aber \(\pi\) ist nur bis auf Einheiten \(\unit=\{\pm1,\pm i\}\)
	und Konjugation bestimmt. Wir fixieren daher eine eindeutige Konvention.
	
	\begin{definition}[Kanonische Wahl \(\pi_p\)]
		Sei \(p\equiv 1\pmod 4\) prim. Wähle die eindeutige Darstellung \(p=a^2+b^2\) mit
		\[
		a>b>0,\qquad a\equiv 1\pmod 2
		\]
		(\(a\) ungerade) und setze
		\[
		\pi_p := a+ib.
		\]
		Falls sowohl \(a+ib\) als auch \(b+ia\) die Bedingungen erfüllen (in der Praxis durch \(a>b\) verhindert), wähle das mit größerem Realteil.
	\end{definition}
	
	\begin{bemerkung}
		Diese Wahl fixiert den Einheitenorbit, weil \(\pi_p\) damit in den ersten Quadranten fällt und zusätzlich \(a\) ungerade ist.
		Alternative kanonische Konventionen (z.\,B.\ \(a\equiv 1\pmod 4\)) sind möglich; für Robustheitstests kann man sie variieren.
	\end{bemerkung}
	
	\begin{lemma}[Kanonische \(\pi_p\) bestimmt \(\arg(\pi_p)\) eindeutig]
		Unter obiger Konvention ist \(\theta_p:=\Arg(\pi_p)\in(0,\pi/2)\) wohldefiniert.
	\end{lemma}
	
	\subsection{Fouriermomente}
	\begin{definition}[Moment-Summen \(S_k(X)\)]
		Für \(k\in\Z\setminus\{0\}\) definieren wir
		\[
		S_k(X):=\sum_{\substack{p\le X\\ p\equiv 1\ (\mathrm{mod}\ 4)}} e^{ik\theta_p},
		\qquad \theta_p:=\Arg(\pi_p).
		\]
		Optional normalisiert man:
		\[
		\widetilde S_k(X):=\frac{1}{\#\{p\le X:\ p\equiv 1(4)\}}\, S_k(X).
		\]
	\end{definition}
	
	\begin{bemerkung}[Bezug zu R3]
		Falls man R3 benutzt (\(z_p=\pi_p^2\) oder \(i\pi_p^2\)), gilt
		\[
		\Arg(z_p) \equiv 2\Arg(\pi_p)\ (\mathrm{mod}\ 2\pi),
		\]
		so dass Fouriermomente der \(\Arg(\pi_p)\) direkt mit Fouriermomenten der \(\Arg(z_p)\) zusammenhängen.
	\end{bemerkung}
	
	\subsection{Nullmodelle für \(S_k(X)\)}
	\begin{definition}[Permutation und Uniform-Nullmodell]
		Zwei einfache Nullmodelle für die Größenordnung von \(|S_k(X)|\) sind:
		\begin{enumerate}
			\item \textbf{Permutation:} Permutiere die Folge \((\theta_{p})\) zufällig über die Primindices und berechne \(S_k\) erneut.
			\item \textbf{Uniformwinkel:} Ersetze \(\theta_p\) durch i.i.d.\ Uniformwinkel auf \((0,\pi/2)\) (oder auf \((-\pi,\pi]\) je nach Konvention).
		\end{enumerate}
		Aus vielen Wiederholungen erhält man Quantile/Bänder für \(|S_k(X)|\).
	\end{definition}
	
	\subsection{Report-Template (Minimalset)}
	\begin{definition}[Fourier-Report (Minimalset)]
		Für \(k\in\{1,2,3,4\}\) und eine Folge \(X_1<X_2<\dots\) soll berichtet werden:
		\begin{enumerate}
			\item \(|\widetilde S_k(X_j)|\) (oder \(|S_k(X_j)|\)) als Funktion von \(X_j\).
			\item Ein Nullmodellband (z.\,B.\ 5\%–95\% Quantile) für \(|\widetilde S_k(X_j)|\).
			\item Robustheitsvergleich: gleiche Auswertung unter Variation der kanonischen \(\pi_p\)-Konvention (mindestens zwei Varianten).
		\end{enumerate}
	\end{definition}
	
	\begin{bemerkung}[Interpretationsregel]
		Ein einzelnes ``auffälliges'' \(k\) bei einem \(X\) ist kein Effekt. Interessant wird es, wenn mehrere \(k\neq 0\)
		gleichzeitig systematisch außerhalb des Nullmodellbandes liegen oder wenn eine strukturierte Oszillation
		(in Skalen \(X\)) auftritt, die unter Robustheitsvariationen erhalten bleibt.
	\end{bemerkung}
	
	% =========================================================
	\section{Zusammenfassung}
	% =========================================================
	Wir haben:
	\begin{enumerate}
		\item die klassischen Teile (Restklasse \(p^2\equiv 1\bmod 12\), Summe-zweier-Quadrate) isoliert,
		\item den Lösungsraum von \(x^2+y^2=p^2\) bis auf Symmetrien und natürliche Kandidaten beschrieben,
		\item die Korridorwahl als Gauge-Fix/Variationsprinzip formalisiert,
		\item und ein reproduzierbares Spiral-Testprotokoll (Unwrap, Score, Nullhypothese, Robustheit) definiert.
	\end{enumerate}
	Damit ist das Gedankenexperiment in eine Form gebracht, die weder überverspricht noch vage bleibt.
	
	% =========================================================
	\begin{thebibliography}{9}
		
		\bibitem{HardyWright}
		G.~H. Hardy und E.~M. Wright,
		\textit{An Introduction to the Theory of Numbers},
		Oxford University Press.
		
		\bibitem{IrelandRosen}
		K.~Ireland und M.~Rosen,
		\textit{A Classical Introduction to Modern Number Theory},
		Springer.
		
		\bibitem{Niven}
		I.~Niven, H.~S. Zuckerman, H.~L. Montgomery,
		\textit{An Introduction to the Theory of Numbers},
		Wiley.
		
	\end{thebibliography}
	
\end{document}